% IEEE conference template — HPML Spring 2026 Final Report
% Team KV: Mamidibathula · Morisetty · Wang · Tian
% Columbia University
\documentclass[conference]{IEEEtran}

\usepackage{amsmath,amssymb,amsfonts}
\usepackage{textcomp}
\usepackage{xcolor}
\usepackage{booktabs}
\usepackage{multirow}
\usepackage{graphicx}
\usepackage{hyperref}
\usepackage{url}
\usepackage{microtype}

\hypersetup{colorlinks=true, linkcolor=blue, citecolor=blue, urlcolor=blue}

% ----------------------------------------------------------------
\begin{document}

\title{KV Cache Optimization for Efficient LLM Inference}

\author{%
\IEEEauthorblockN{%
  Rishika Mamidibathula\textsuperscript{\dag},
  Suhas Morisetty\textsuperscript{*},
  Ziheng Wang\textsuperscript{*},
  Tony Tian\textsuperscript{*}}
\IEEEauthorblockA{%
  \textsuperscript{*}Department of Computer Science,
  \textsuperscript{\dag}Department of Data Science,
  Columbia University\\
  COMS E6998: High Performance Machine Learning, Spring 2026\\
  \{rm4318, vm2825, zw3153, jt3640\}@columbia.edu}
}

\maketitle

% ----------------------------------------------------------------
\begin{abstract}
Autoregressive large-language-model (LLM) decoding is fundamentally
bounded by key-value (KV) cache memory bandwidth.  At batch size 32
and a 4K context, Llama-2-7B alone requires roughly 65\,GB for the KV
cache, nearly saturating an H100 80\,GB GPU.  We present a controlled
benchmark of five KV optimization techniques drawn from four
compression families---per-token quantization (KIVI 2/4-bit), dynamic
sparse selection (TopK / TokenSelect), one-shot eviction (SnapKV), and
latent projection (TransMLA)---together with a quantized-sparse hybrid
(KIVI$\times$TopK), all evaluated under a single unified harness that
fixes the model, hardware, decode loop, prompts, and random seeds
across runs.  KIVI 4-bit preserves baseline LongBench quality with a
$1.76\times$ decode-throughput gain at batch 32, while KIVI 2-bit
trades $-6.2\%$ overall quality for a $1.93\times$ gain at batch 32 and
lifts the operational batch ceiling from 32 to 128 (where baseline
OOMs).  SnapKV marginally improves overall quality at no per-step
decode cost, though its one-shot eviction heuristic may not generalise
uniformly beyond controlled benchmarks.  TopK $K\!=\!1024$ matches
baseline quality and gains substantially on retrieval-heavy tasks.
MLA yields a consistent $3.9\times$ KV reduction even at batch 1, but
the available checkpoint regresses on output quality and requires
retraining at longer context to be production-ready.  The hybrid
method stacks the memory savings of both compression axes, achieving a
50\% reduction at 32K context; however, decode throughput degrades due
to a non-optimal kernel implementation that does not yet fuse scoring,
selection, and attention into a single pass---suggesting the latency
gap is an engineering artifact rather than a fundamental limitation.
Overall, each method occupies a distinct point in the
quality--memory--latency space, and the right choice depends on
whether batch size, context length, or output fidelity is the binding
constraint.
\end{abstract}

\begin{IEEEkeywords}
KV cache, LLM inference, quantization, sparse attention, token eviction,
multi-head latent attention, throughput, memory efficiency
\end{IEEEkeywords}

% ================================================================
\section{Introduction}
% ================================================================

LLM inference proceeds in two phases.  Prefill processes the prompt in
a single forward pass; decode then emits one token per step.  Each
decode step rereads every cached key and value tensor from HBM, making
decode strongly memory-bandwidth bound rather than compute-bound.  The
cost grows linearly with sequence length, batch size, and depth.  For
Llama-2-7B \cite{llama2} (32 layers, 32 heads of dimension 128, FP16) at
4096 tokens and batch size 32, the cache occupies
$2\times32\times32\times128\times4096\times32\times2 \approx 65$\,GB
\cite{survey}---leaving little room for activations on an 80\,GB GPU
and pushing larger configurations into out-of-memory territory.

Research on KV reduction has converged on four principal compression
families.  \emph{Quantization} reduces bits per element while keeping
all tokens \cite{kivi,kvquant}.  \emph{Sparse selection} attends to a
dynamically chosen subset of tokens at each decode step
\cite{tokenselect,rocketkv}.  \emph{Eviction} permanently discards
low-importance tokens after prefill \cite{h2o,snapkv,adakv}.
\emph{Latent projection} stores a compressed low-rank representation
instead of full K/V tensors \cite{deepseekv2,transmla}.  Cross-family
comparison is difficult in practice because published evaluations
rarely share the same model, hardware, or decode loop.  This work
addresses that gap with a unified harness whose only variable is the KV
policy.

Our contributions are: (i) a single-codebase benchmark of five KV
optimization techniques and a quantized-sparse hybrid on Llama-2-7B;
 (ii) empirical evidence that KIVI 4-bit and SnapKV 0.4
deliver lossless or marginally better LongBench quality; and (iii) full
reproducibility via Modal cloud scripts and a public W\&B dashboard at
\url{https://wandb.ai/rm4318-columbia-university/kv-cache-hpml}.

% ================================================================
\section{Literature Review}
% ================================================================

\textbf{KV Cache Compression (Quantization).}
Low-bit quantization reduces bytes per stored token while preserving
full token coverage.  KIVI \cite{kivi} applies asymmetric quantization
with per-channel scales for keys and per-token grouped quantization
(group size 32) for values at 2--4 bit precision, retaining a
recent-token window in FP16 for short-range fidelity; quantization is
applied strictly at inference time without retraining or weight
modification.  At 2-bit it achieves roughly $8\times$ cache
compression, at 4-bit roughly $4\times$.  KVQuant \cite{kvquant}
extends this with non-uniform codebooks and per-channel smoothing,
reaching near-lossless compression at 1-bit at the cost of greater
overhead.

\textbf{KV Cache Low-Rank Approximation.}
Low-rank methods reduce per-token dimensionality rather than token
count by projecting cached KV tensors into a lower-dimensional space.
xKV \cite{xkv} applies layer-group-wise SVD directly to cached KV
tensors, stores the projected coefficients, and reconstructs
approximate K/V at attention time---remaining strictly within
inference scope by operating on the cache rather than model weights.
DeepSeek-V2 \cite{deepseekv2} introduced Multi-Head Latent Attention
(MLA), which stores a compressed latent vector $c_t$ per token and
expands to full K/V via learned up-projections, with Attention
Absorption amortizing the projection cost.  TransMLA \cite{transmla}
retrofits this latent structure onto an existing Llama-2 checkpoint,
enabling evaluation of low-rank projection without full retraining.

\textbf{KV Cache Retention Management (Eviction).}
Eviction methods permanently discard low-importance tokens to reduce
stored token count while preserving high-utility context.  H2O
\cite{h2o} identifies ``heavy-hitter'' tokens by cumulative attention
mass and retains a recent window alongside these high-importance
tokens.  SnapKV \cite{snapkv} performs a single one-shot eviction
after prefill by aggregating attention weights from a small
observation window and retaining the top-scored tokens per head,
producing a fixed-size compressed cache with no per-step overhead.
Ada-KV \cite{adakv} extends this by allocating the eviction budget
adaptively per head rather than uniformly, reflecting varying
information density across attention heads.

\textbf{Token-Level KV-Centric Scheduling (Sparse Selection).}
Unlike eviction, token-level scheduling reduces effective attention
work without permanently removing KV states---masked tokens remain
stored and may be reactivated.  At each decode step, a lightweight
importance score (approximate $\mathbf{q}\cdot\mathbf{k}$) is
computed and only the top-$K$ tokens participate in attention, with
the remainder masked.  TokenSelect \cite{tokenselect} implements this
with a fused Triton scoring kernel and a cosine-similarity reuse cache
to avoid re-scoring when consecutive queries are similar.  RefreshKV
\cite{refreshkv} augments this with periodic full-context refresh to
recover information from previously masked tokens.  RocketKV
\cite{rocketkv} adds a two-stage compression pipeline to further
reduce prefill and decode costs in long-context settings.

% ================================================================
\section{Methodology}
% ================================================================

\subsection{Model and Hardware}

All experiments use \texttt{meta-llama/Llama-2-7b-chat-hf} (32 layers,
32 heads, head dim 128, FP16 weights) on a single NVIDIA H100 80\,GB
HBM3 GPU provided by Modal cloud compute.  The runtime stack is pinned:
PyTorch 2.4.1, HuggingFace Transformers 4.44.2, CUDA 12.1, Triton 3.0,
flash-attn 2.5.  The MLA experiments additionally use
\texttt{botsxc/llama2-7b-chat-mla-2048-8} \cite{transmla} (2048-dim
latent, 256-token calibration context, no further finetuning).

\subsection{Datasets and Metrics}

\textbf{Quality.}  LongBench \cite{longbench}, six tasks $\times$ 20
examples = 120 examples per method: \texttt{qasper} (scientific QA, F1),
\texttt{multifieldqa\_en} (multi-document QA, F1), \texttt{triviaqa}
(open-domain QA, F1), \texttt{2wikimqa} (multi-hop QA, F1),
\texttt{multi\_news} (summarization, ROUGE-L), \texttt{lcc} (code
completion, edit similarity).  Inputs are middle-truncated to 4096
tokens following the KIVI paper convention.

\textbf{Throughput.}  Synthetic fixed-size batches with prefill 1024
and generation 512 tokens, swept over batch sizes 1--256.  Five batches
per configuration; the first is dropped as warmup; the mean of the
remaining four is reported.

\textbf{Long context.}  Prefill 32{,}768 tokens, generation 512, batch
size 1, used for TopK-Flash and the KIVI$\times$TopK hybrid.

\subsection{Profiling Tools}

Memory: \texttt{torch.cuda.max\_memory\_allocated()} reset before each
batch; this measures peak, not steady-state.  Timing:
\texttt{torch.cuda.synchronize()} + \texttt{time.perf\_counter()}
brackets prefill (TTFT, time to first token) and the decode loop
(throughput in tokens per second over steps $2$ to $G$).
Operator-level traces are captured with \texttt{torch.profiler} and
exported to Chrome trace.  All runs are logged to Weights \& Biases.

\subsection{Method Implementations}

All methods implement a shared \texttt{MethodWrapper} interface in
\texttt{methods/} and run through an identical generate loop in
\texttt{benchmark/runner.py}.

\textbf{Baseline (FP16):} dense attention with no KV modification.

\textbf{KIVI \cite{kivi}.}  Keys are quantized per-channel to 2 or 4
bits (grouped across the sequence dimension per head); values are
quantized per-token with group size 32.  The last 128 tokens are kept
in FP16 (the residual buffer).  Quantization fuses INT4 packing in a
Triton kernel (\texttt{triton\_quantize\_and\_pack}); GEMV dequantizes
on-the-fly via \texttt{cuda\_bmm\_fA\_qB\_outer}.  The cache is stored
as a 9-tuple per layer (quantized blocks + scales + zeros for K and V
plus the FP16 residual).

\textbf{TopK / TokenSelect \cite{tokenselect}.}  The full FP16 cache is
retained; only the attention \emph{scope} is sparse.  Each decode step
divides the cache into three non-overlapping regions: a fixed sink
window (first 128 tokens), a dynamically selected top-$K$ middle block,
and a fixed local window (last 512 tokens).  Scoring is
per-layer-independent: a Triton kernel (\texttt{\_topk\_score\_kernel})
fuses $\mathbf{Q}\mathbf{K}^\top / \sqrt{d}$, head-wise softmax (soft
voting per paper Eq.~7), and a cross-head sum into a single GPU launch,
then \texttt{torch.topk} selects the top-$K$ middle indices.  An
optional max-pool smoothing (kernel size configurable) may be applied
before top-$K$ to cluster neighbouring tokens.  A per-layer
cosine-similarity cache (threshold 0.9) reuses the previous step's
indices when consecutive query fingerprints are sufficiently similar,
avoiding redundant scoring.  Because the selected subset is shorter than
the original cache, a \texttt{true\_seq\_length} counter (= full
accumulated cache length) is passed as \texttt{position\_ids} at each
decode step to keep RoPE indices correct.  At long context (32K) a
paged variant, \textbf{TopK-Flash}, uses a 256-token paged KV pool with
\texttt{flash\_attn\_with\_kvcache} for zero-copy selective attention,
eliminating the ${\sim}52$\,GB contiguous allocation that standard
flash-attn requires.

\textbf{SnapKV \cite{snapkv}.}  A single one-shot eviction is performed
immediately after prefill; decode then runs on the fixed compressed
cache at zero per-step overhead.  The algorithm has three stages.
\emph{Vote}: attention weights from the last $W\!=\!32$ query positions
(the \emph{observation window}) are summed over the prefix per head,
yielding an importance score of shape
$(\text{batch},\,\text{heads},\,\text{prefix\_len})$.  Selection is
done \emph{per head}---averaging across heads would erase per-head
specialisation.  \emph{Pool}: a 1D max-pool with kernel size 7 and
same-padding is applied per head across the prefix axis \emph{before}
top-$k$ selection to cluster neighbouring tokens and prevent fragmented
picks (matching paper Section 4.3).  \emph{Snap}: the top-$\lfloor r
\times S \rfloor$ prefix positions per head are gathered (budget ratio
$r\!=\!0.4$, sequence length $S$); the full observation window is then
concatenated back unconditionally, so no special sink bookkeeping is
needed for it.  Optionally, a configurable \texttt{sink\_size} prepends
the first few tokens (attention sinks) to the retained set.  After
eviction the DynamicCache is shorter than the original prompt; a
\texttt{true\_seq\_length} counter ($=$ original prefill length $+$
decode steps elapsed) is threaded to each decode step so that RoPE
position indices remain accurate despite the shortened cache.

\textbf{MLA / TransMLA \cite{transmla}.}  Each Llama-2 attention layer
is replaced with a DeepSeek-V3-style MLA module.  A down-projection
$\mathbf{W}_{\text{kv}}^A$ maps hidden states to a low-rank latent
$c_t^{\text{raw}} \in \mathbb{R}^{d_c}$ and a RoPE component
$k_t^R \in \mathbb{R}^{d_r}$ (split from a single linear),
where $d_c\!=\!2048$ (\texttt{kv\_lora\_rank}) and $d_r\!=\!64$
(\texttt{qk\_rope\_head\_dim}) for the evaluated checkpoint.  An
up-projection $\mathbf{W}_{\text{kv}}^B$ expands $c_t$ back to full
$(\mathbf{K}_{\text{nope}},\mathbf{V})$ tensors at attention time.

Our implementation uses a custom \texttt{MLALatentCache} that stores
\emph{only} the normalised latent $\tilde{c}_t =
\text{LayerNorm}(c_t^{\text{raw}})$ and the RoPE-applied rotation
$\hat{k}_t^R$ per token per layer, rather than the fully expanded K/V
pair.  At each decode step, $\mathbf{W}_{\text{kv}}^B(\tilde{c}_{1:t})$
is evaluated over all accumulated latents to reconstruct K and V
on-the-fly---an $O(t)$ compute cost that trades FLOP for memory.  This
avoids the pitfall of the na\"{\i}ve path (storing expanded K/V in a
\texttt{DynamicCache}), which would be \emph{larger} than the baseline
GQA cache given the expanded head dimensions.  Per-token cache
footprint for the latent path:
$(2048 + 64) \times 2$\,B/layer $\times 32$ layers
$\approx 132$\,KB/token, versus $\approx 512$\,KB/token for the
FP16 baseline ($32\,\text{heads}\times 128\,d_k \times 2 \times 2$\,B
$\times 32$ layers), yielding the measured $3.9\times$ KV reduction
(e.g., $1207 \rightarrow 311$\,MB at prefill 2048, BS=1).

\textbf{KIVI$\times$TopK Hybrid.}  Prefill uses
\texttt{F.scaled\_dot\_product\_attention} (flash-attn backed,
$O(T)$ memory) to avoid 32K OOM.  Decode scores all tokens via
\texttt{cuda\_bmm\_fA\_qB\_outer} over INT4-packed K, selects
$n_\text{sink}+K+n_\text{local}$ positions, and gathers the
corresponding INT4 V rows into a contiguous buffer fed directly into
\texttt{cuda\_bmm\_fA\_qB\_outer}---no intermediate FP16 V is
materialized.

% ================================================================
\section{Experimental Results}
% ================================================================

\subsection{LongBench Quality}

Table~\ref{tab:quality} reports per-task and overall LongBench scores.
KIVI 4-bit and TopK $K\!=\!1024$ both achieve 0.292 ($+0.001$ absolute,
$+0.3\%$ relative over baseline 0.291)---within evaluation noise for
20 examples per task but consistent across runs.  SnapKV 0.4 reaches
0.295 ($+0.004$ absolute, $+1.4\%$ relative).

\begin{table}[!t]
\centering
\caption{LongBench quality (F1 / ROUGE-L / edit-sim), 6 tasks $\times$ 20 ex.
Tasks: (1) qasper, (2) mfqa, (3) triviaqa, (4) 2wikimqa,
(5) mnews, (6) lcc.}
\label{tab:quality}
\setlength{\tabcolsep}{3.5pt}
\begin{tabular}{@{}lrrrrrrr@{}}
\toprule
\textbf{Method} & \textbf{(1)} & \textbf{(2)} & \textbf{(3)} & \textbf{(4)} & \textbf{(5)} & \textbf{(6)} & \textbf{Avg} \\
\midrule
Baseline (FP16)       & .172 & .380 & .644 & .083 & .181 & .287 & .291 \\
KIVI 4-bit            & .194 & .360 & .643 & .083 & .188 & .286 & \textbf{.292} \\
KIVI 2-bit            & .133 & .382 & .608 & .050 & .191 & .277 & .273 \\
TopK $K$=2048         & .145 & .329 & .731 & .075 & .178 & .282 & .290 \\
TopK $K$=1024         & .153 & .281 & .773 & .089 & .183 & .273 & \textbf{.292} \\
TopK $K$=512          & .145 & .237 & .751 & .050 & .178 & .281 & .274 \\
SnapKV 0.4            & .153 & .387 & .687 & .083 & .181 & .277 & \textbf{.295} \\
\midrule
Baseline (2K trunc)$^{\ddagger}$ & .124 & .298 & .375 & .089 & .180 & .362 & .238 \\
MLA (2K trunc)$^{\ddagger}$      & .036 & .101 & .064 & .043 & .103 & .140 & .081 \\
\bottomrule
\end{tabular}
\vspace{1pt}
\footnotesize{$^{\ddagger}$MLA evaluation uses 2048-token middle truncation rather
than 4096 to match the checkpoint's calibration context; the matched
2K-truncation baseline (0.238) is the appropriate reference for the
MLA row, yielding $-66\%$ relative.  Numbers in this block are not
directly comparable to the 4K-truncation rows above.}
\end{table}

Per-task patterns reveal two opposite breakdown mechanisms when
compression is too aggressive.  KIVI 2-bit (0.273) and TopK
$K\!=\!512$ (0.274) end up at essentially the same overall score, but
for different reasons.  KIVI's quantization noise hurts qasper
($-22\%$) and 2wikimqa ($-40\%$), tasks that require precise recall of
sparse mid-context facts.  TopK $K\!=\!512$ hurts
\texttt{multifieldqa\_en} ($-38\%$) because answer spans deeper than
$n_\text{sink} + n_\text{local}$ are pruned.  Conversely, TopK beats
baseline on TriviaQA across all $K$ ($+13$\,pp to $+20$\,pp absolute):
the answer appears in multiple paraphrases, so sparse selection drops
distracting copies and concentrates attention.  SnapKV's one-shot vote
actively aligns selection with answer-bearing positions, boosting
triviaqa ($+4.3$\,pp) and multifieldqa ($+0.7$\,pp) while only hurting
qasper ($-2.0$\,pp) where evidence spans fall outside the 32-token
observation window.  The MLA row, evaluated against its matched
2K-truncation baseline (0.238), drops to 0.081---a $-66\%$ relative
regression---with the largest collapses on retrieval-heavy tasks
(triviaqa $-83\%$, multifieldqa $-66\%$); see \S\ref{sec:limitations}
for the underlying cause.

\subsection{Throughput and Memory}

Table~\ref{tab:throughput} shows the KIVI throughput sweep at prefill
1024, generation 512, for both 4-bit and 2-bit precision.

\begin{table}[!t]
\centering
\caption{Decode throughput (tok/s) and peak GPU memory vs.\ batch size.
Prefill 1024, gen 512.  ``--'' = OOM. Speedup vs.\ FP16 baseline.}
\label{tab:throughput}
\setlength{\tabcolsep}{2.5pt}
\footnotesize
\begin{tabular}{@{}rrrrrrrrrr@{}}
\toprule
 & \multicolumn{5}{c}{\textbf{Decode tok/s}} & \multicolumn{4}{c}{\textbf{Peak mem (GB)}} \\
\cmidrule(lr){2-6} \cmidrule(lr){7-10}
\textbf{BS} & \textbf{Base} & \textbf{4b} & \textbf{Spd} & \textbf{2b} & \textbf{Spd} & \textbf{Base} & \textbf{4b} & \textbf{2b} & \textbf{MLA} \\
\midrule
1   &  70 &   29 & 0.41$\times$ &   30 & 0.43$\times$ & 15.1 & 14.3 & 14.0 & 13.8 \\
2   & 135 &   58 & 0.43$\times$ &   61 & 0.45$\times$ & 16.7 & 14.8 & 14.4 &  --   \\
4   & 242 &  109 & 0.45$\times$ &  123 & 0.51$\times$ & 20.0 & 16.0 & 15.3 &  --   \\
8   & 351 &  216 & 0.62$\times$ &  243 & 0.69$\times$ & 26.4 & 18.4 & 17.1 &  --   \\
16  & 438 &  404 & 0.92$\times$ &  480 & 1.10$\times$ & 39.3 & 23.7 & 20.7 &  --   \\
32  & 497 &  873 & 1.76$\times$ &  957 & 1.93$\times$ & 65.0 & 33.6 & 27.8 &  --   \\
64  &  -- & 1186 & $\infty$     & 1295 & $\infty$     &  --  & 53.6 & 42.1 &  --   \\
128 &  -- &  --  & ---          & 1517 & $\infty$     &  --  &  --  & 70.7 &  --   \\
\bottomrule
\end{tabular}
\vspace{2pt}

\footnotesize{Columns ``4b'' and ``2b'' are KIVI 4-bit and 2-bit
respectively.  KIVI 4-bit OOMs at BS=128 because its per-token
footprint is twice that of 2-bit; KIVI 2-bit is the only configuration
that fits BS=128 on a single H100.}
\end{table}

\textbf{Crossover at BS=16.}  KIVI 4-bit decode throughput is
$0.4$--$0.6\times$ baseline at BS 1--8, reaches near-parity at BS=16
($404$ vs.\ $438$\,tok/s, $0.92\times$), and crosses to $1.76\times$
at BS=32; KIVI 2-bit follows the same shape but crosses earlier (at
BS=16, $1.10\times$) and reaches $1.93\times$ at BS=32.  The mechanism
is straightforward: at small batch the decode
$\mathbf{Q}\mathbf{K}^\top$ is a GEMV, and INT4 GEMV does not map to
H100 tensor-core units---it is roughly $2.5\times$ slower than FP16
GEMV.  At large batch the same operation becomes a near-GEMM where the
HBM bandwidth saving (the 2-bit cache is $2.3\times$ smaller than
baseline at BS=32: 27.8 vs.\ 65.0\,GB; 4-bit is $1.9\times$ smaller at
33.6\,GB) outweighs dequantization cost.

\textbf{Capacity extension.}  Baseline OOMs at BS=64; KIVI 4-bit
reaches BS=64 (53.6\,GB, 1186\,tok/s) and KIVI 2-bit reaches BS=128
(70.7\,GB, 1517\,tok/s)---about $3\times$ the maximum throughput ever
achievable with the FP16 baseline on this GPU.  This is the strongest
practical argument for KIVI: it raises the per-GPU concurrency
ceiling, not just speed at fixed batch.  KIVI 4-bit is the
quality-preserving operating point (0.292 vs.\ baseline 0.291); KIVI
2-bit pays $-6.2\%$ overall quality (0.273) for an additional batch
doubling and modestly higher throughput.

\textbf{TTFT overhead.}  KIVI's TTFT is 18--66\% worse across batch
sizes (998 vs.\ 1659\,ms at BS=32 for 2-bit; 4-bit is essentially
identical) because the full KV cache is quantized and packed at the
end of prefill---a cost proportional to $B \times S$.  For
latency-sensitive interactive workloads this is a real penalty; for
throughput-oriented offline inference it is amortized over hundreds of
decode steps.

\textbf{MLA at B=1.}  At gen=256, MLA achieves 43--44\,tok/s vs.\
34\,tok/s baseline ($1.26$--$1.32\times$) across prefill lengths
512--2048.  Even single-request decode is partially HBM-bandwidth
bound: the $3.9\times$ KV reduction (e.g., 1207 \(\rightarrow\) 311\,MB
at $p\!=\!2048$) reduces per-step HBM traffic proportionally.  The
speedup narrows at gen=512 because the MLA up-projection cost grows
with output length.

\subsection{Long-Context Results (32K)}

\begin{table}[!t]
\centering
\caption{Long-context results.  Prefill 32{,}768 tokens, gen 512, BS=1.}
\label{tab:longctx}
\begin{tabular}{@{}lrrrr@{}}
\toprule
\textbf{Method} & \textbf{TTFT (ms)} & \textbf{Spdup} & \textbf{Peak mem} & \textbf{Saving} \\
\midrule
Baseline (FP16)     & 2627 & 1.00$\times$ & 52.1\,GB & --- \\
TopK-Flash          & 1761 & 1.49$\times$ & 42.1\,GB & $-19\%$ \\
SnapKV 0.4          &  --- & ---          & 26.0\,GB$^{\dagger}$ & $-50\%$ \\
KIVI$\times$TopK 4b &  --- & 0.61$\times$ & 25.8\,GB & $-50\%$ \\
\bottomrule
\end{tabular}
\vspace{1pt}
\footnotesize{Baseline OOMs at 64K; TopK-Flash succeeds at 65.8\,GB.
$^{\dagger}$SnapKV peak memory at 32K is interpolated from
short-context measurements; 32K SnapKV throughput not directly
benchmarked.}
\end{table}

\textbf{TopK-Flash.}  At 32K context the 256-token paged KV pool
eliminates the 52\,GB contiguous allocation that standard flash-attn
requires, reducing TTFT by 33\% and peak memory by 19\%.  The win is
not from sparsity but from allocation and memory-access-pattern
improvement.  Baseline OOMs at 64K; TopK-Flash succeeds at 65.8\,GB.

\textbf{KIVI$\times$TopK hybrid.}  Combining INT4 storage with sparse
attention cuts peak memory 50\% (52.1$\rightarrow$25.8\,GB) but decode
throughput drops 39\% (21.0$\rightarrow$12.9\,tok/s).  The slowdown has
two attributable sources: (a) $\mathbf{Q}\mathbf{K}^\top$ must still
score \emph{all} 32K tokens every step before top-$K$ can be
applied---the scoring cost is identical to dense baseline, plus the
overhead of topk + sort + gather on top; (b) at batch size 1 and 32K
context the INT4 V GEMV inherits the same dequantization penalty as
standalone KIVI, where tensor-core utilisation is low and the
dequantization overhead is not amortised.  LongBench quality remains
at 0.297, within noise of 0.291, confirming the compression is
semantically lossless.

Importantly, we believe the throughput regression is \emph{not}
fundamental to the hybrid strategy but rather an artifact of the
current kernel implementation.  The scoring kernel reads the full INT4
K cache sequentially before any selection, and the subsequent gather
plus INT4 V attention are not fused---each is a separate kernel launch
with its own HBM round-trip.  A tighter implementation that fuses
scoring, index selection, and INT4 V attention into a single paged
kernel (analogous to what FlashInfer or FlashAttention-3 do for dense
attention) could in principle recover the latency by eliminating
redundant memory traffic.  Similarly, a hierarchical or lazy scoring
scheme---e.g., scoring only a coarse page-level summary first and
expanding to token-level only for shortlisted pages---would reduce the
$O(S)$ scoring cost toward $O(K)$.  We therefore view the hybrid as a
promising direction that merits further kernel-level engineering: the
memory savings from orthogonal INT4 and sparse-selection compression
are real and additive, and achieving the best of both worlds
(throughput \emph{and} memory) appears within reach with a
purpose-built fused kernel.

% ================================================================
\section{Discussion}
% ================================================================

\subsection{Deployment Recommendations}

Table~\ref{tab:tradeoff} summarizes the operating point of each method.

\begin{table}[!t]
\centering
\caption{Deployment fit summary.}
\label{tab:tradeoff}
\setlength{\tabcolsep}{3pt}
\begin{tabular}{@{}llll@{}}
\toprule
\textbf{Method} & \textbf{KV ratio} & \textbf{Quality} & \textbf{Best regime} \\
\midrule
KIVI 4-bit  & $4\times$  & Lossless (.292)        & Batch serving $\geq$ BS=16 \\
KIVI 2-bit  & $8\times$  & $-6.2\%$ (.273)        & Max-batch, accept noise \\
TopK $K$=1024 & $1\times$$^{\dagger}$ & Lossless (.292) & Single-req, long context \\
SnapKV 0.4  & $2.5\times$ & $+1.4\%$ (.295)       & Any; quality-neutral \\
MLA (B=1)   & $3.9\times$ & $-66\%$* needs retrain & After finetuning at $\geq$4K \\
\bottomrule
\end{tabular}
\vspace{1pt}
\footnotesize{$^{\dagger}$TopK keeps full FP16 cache; only attention scope is sparse.  *Measured against matched 2K-truncation baseline (0.238).}
\end{table}

\textbf{KIVI 4-bit} is the best practical default for high-throughput
batch inference: it costs nothing in quality, raises max batch from 32
to 64 on a single H100 (and 128 if 2-bit is acceptable), and
$1.76\times$es throughput at BS=32.  The cost is TTFT latency
($+18$--$66\%$), which matters for interactive applications.  KIVI
2-bit is the right choice when maximum concurrency dominates and a
$-6.2\%$ quality cost is tolerable.

\textbf{SnapKV 0.4} is the right choice when quality must be preserved
and per-step overhead must be zero.  Its $2.5\times$ KV reduction has
no decode-time cost because the compressed cache is built once after
prefill.  That said, our benchmark uses only 20 LongBench examples per
task with inputs truncated to 4096 tokens, and the one-shot eviction
heuristic (voting over a 32-token observation window) may not generalise
equally well to all real-world distributions.  Tasks where critical
evidence is spread uniformly across the prompt, or where the most
relevant tokens do not align with the final observation window, could
see larger quality degradation than our controlled evaluation suggests.
More deliberation is needed---across longer contexts, more diverse task
types, and multi-turn settings---before SnapKV can be recommended as a
broadly quality-neutral drop-in.

\textbf{TopK $K\!=\!1024$} is suited to single-request, long-context
use ($>32$K) where the paged variant (TopK-Flash) provides real TTFT
and memory gains.  At short context and low batch, scoring overhead
dominates the saved attention work.

\textbf{MLA} shows the theoretical ceiling of latent projection:
$3.9\times$ KV reduction and $1.3\times$ decode speedup are repeatable
benefits even at B=1, but production-quality output requires a
checkpoint trained at longer context with higher latent rank.

\subsection{Kernel-Level Observations}

\textbf{INT4 GEMV inefficiency.}  \texttt{cuda\_bmm\_fA\_qB\_outer}
performs on-the-fly INT4 $\rightarrow$ FP16 dequantization per GEMV.
H100 tensor cores target INT8/FP16 GEMM; INT4 GEMV falls back to scalar
or vector-INT4 paths.  Bandwidth saving only exceeds this overhead when
decode approaches GEMM at BS$\geq$16.

\textbf{KIVI residual buffer hot path.}  At short generation lengths
($<\!\texttt{residual\_length}\!\times\!n_\text{layers}$), nearly all
attention concentrates in the FP16 residual; KIVI provides minimal
storage benefit in this regime.  This is also why KIVI TTFT is roughly
$3\times$ slower: the full prefill quantization is paid upfront.

\textbf{Triton fused scoring.}  \texttt{\_topk\_score\_kernel} fuses
$\mathbf{Q}\mathbf{K}^\top$ + softmax + head-sum into a single launch,
reducing Python-dispatch overhead $2$--$3\times$ vs.\ sequential
PyTorch ops.  At 1K--4K sequence lengths the critical path remains HBM
bandwidth, not compute or launch latency.

\textbf{Hybrid stays in INT4.}  The hybrid keeps data in INT4 across
the gather $\rightarrow$ attention boundary by passing the gathered
packed tensor directly into \texttt{cuda\_bmm\_fA\_qB\_outer},
avoiding an FP16 V materialization that would have cost an extra
allocation and HBM round-trip per layer per step.

\subsection{Limitations}
\label{sec:limitations}

\emph{Single architecture.}  Results are specific to Llama-2-7B with
multi-head attention (32 heads, 32 KV heads).  Models with grouped-query
attention (Llama-3, Mistral) have fewer KV heads, which changes the
compression profile.

\emph{Small evaluation set.}  LongBench results use 20 examples per
task.  Some per-task deltas (e.g., qasper KIVI 2-bit $-22\%$) may
partially reflect small-$n$ variance, although the direction is
consistent across tasks.

\emph{MLA checkpoint and truncation.}  The available TransMLA
checkpoint was calibrated at 256-token context with rank-8 latent
compression and without task finetuning---insufficient capacity for
4K-token retrieval-heavy LongBench tasks.  We accordingly evaluated
MLA against a matched 2K-truncation baseline (0.238) rather than the
4K-truncation baseline used for the other methods (0.291); the
$-66\%$ quality gap reported in Table~\ref{tab:tradeoff} reflects
this matched comparison.  A properly trained checkpoint at higher
rank (16--32) and longer context (4--8K) would likely recover most of
this gap, as the underlying latent-projection mechanism produces a
clean $3.9\times$ KV reduction and $1.3\times$ decode speedup.

\emph{Hybrid scoring cost.}  The 50\% memory win comes with a 39\%
decode penalty because scoring must precede selection.  Hierarchical
or lazy-scoring approaches could decouple them.

% ================================================================
\section{Conclusion}
% ================================================================

We conducted a controlled benchmark of five KV-cache optimization
methods spanning four compression families on Llama-2-7B.  KIVI 4-bit
provides lossless quality with $4\times$ KV compression and
$1.76\times$ throughput at BS=32---our recommended default for
batch-oriented serving when quality must be preserved; KIVI 2-bit
extends the operational batch ceiling from 32 to 128 at a $-6.2\%$
quality cost.  SnapKV marginally improves LongBench overall score with
zero runtime overhead.  TopK $K\!=\!1024$ matches baseline quality and
gains substantially on retrieval-heavy tasks.  MLA delivers a
consistent $3.9\times$ KV reduction and $1.3\times$ decode speedup at
B=1 but requires a properly trained checkpoint for production quality.
The KIVI$\times$TopK hybrid stacks memory savings of both axes but
pays a decode penalty until scoring cost can be reduced.  No single
method dominates all axes; practitioners should select based on
whether batch size, context length, or output fidelity is the binding
constraint.

Future work should address: (i)~batched decode for TopK to amortize
the per-step scoring cost and produce a KIVI-style crossover at
BS$\geq$8--16; (ii)~MLA retraining at 4K--8K context with rank 16--32
to preserve $3.9\times$ compression while recovering QA quality;
(iii)~hierarchical or lazy scoring in the KIVI$\times$TopK hybrid to
decouple the memory win from the decode penalty; and (iv)~cross-method
composition, e.g., SnapKV eviction followed by KIVI quantization on
the retained tokens to stack savings from orthogonal axes.

\appendix
\section*{AI Use Disclosure}

Per the HPML AI Use Policy.

\textbf{Tools used:} Claude (Anthropic) and GitHub Copilot.

\textbf{Specific purpose:} Debugging CUDA OOM errors during KIVI
kernel integration; clarifying Triton kernel semantics (block-pointer
arithmetic, mask construction, autotune configuration); polishing
prose in sections drafted by team members.

\textbf{Sections affected:} \texttt{methods/kivi\_kernels/} integration
debugging, \texttt{methods/topk\_selection.py} kernel review;
\S{}I--\S{}VI prose polish on text first written by team members.  No
analytical reasoning, performance interpretation, or experimental
result was generated by an AI tool.

\textbf{How we verified correctness:} All experimental numbers in this
report were produced by running our own Modal scripts on H100
containers and logged to W\&B; profiler-trace interpretations were
cross-checked against the raw \texttt{torch.profiler} traces in
\texttt{results/}; AI-suggested code was rewritten in our own words
and re-validated to produce identical outputs to a hand-written
reference before being committed.

By submitting this project, the team confirms that the analysis,
interpretations, and conclusions are our own, and that any AI
assistance is fully disclosed above.  The same disclosure block
appears in the project README.

% ================================================================
\bibliographystyle{IEEEtran}

\begin{thebibliography}{14}

\bibitem{survey}
J.~Jiang, P.~Yang, R.~Zhang, and F.~Liu,
``Towards efficient large language model serving: A survey on
system-aware KV cache optimization,''
\textit{TechRxiv}, 2026.

\bibitem{llama2}
H.~Touvron \textit{et al.},
``Llama 2: Open foundation and fine-tuned chat models,''
\textit{arXiv:2307.09288}, 2023.

\bibitem{kivi}
Z.~Liu \textit{et al.},
``KIVI: A tuning-free asymmetric 2-bit quantization for KV cache,''
in \textit{Proc.\ ICML}, 2024.

\bibitem{kvquant}
C.~Hooper \textit{et al.},
``KVQuant: Towards 10 million context length LLM inference with KV
cache quantization,''
in \textit{NeurIPS}, 2024.

\bibitem{tokenselect}
W.~Wu \textit{et al.},
``TokenSelect: Efficient long-context inference via dynamic token-level
KV cache selection,''
in \textit{Proc.\ EMNLP}, 2025.

\bibitem{rocketkv}
P.~Behnam \textit{et al.},
``RocketKV: Accelerating long-context LLM inference via two-stage KV
cache compression,''
in \textit{Proc.\ ICML}, 2025.

\bibitem{refreshkv}
F.~Xu, T.~Goyal, and E.~Choi,
``RefreshKV: Updating small KV cache during long-form generation,''
in \textit{Proc.\ ACL}, 2025.

\bibitem{h2o}
Z.~Zhang \textit{et al.},
``H2O: Heavy-hitter oracle for efficient generative inference of large
language models,''
in \textit{NeurIPS}, 2023.

\bibitem{snapkv}
Y.~Li \textit{et al.},
``SnapKV: LLM knows what you are looking for before generation,''
in \textit{NeurIPS}, 2024.

\bibitem{adakv}
Y.~Feng \textit{et al.},
``Ada-KV: Optimizing KV cache eviction by adaptive budget allocation,''
in \textit{NeurIPS}, 2025.

\bibitem{deepseekv2}
DeepSeek-AI,
``DeepSeek-V2: A strong, economical, and efficient mixture-of-experts
language model,''
\textit{arXiv:2405.04434}, 2024.

\bibitem{transmla}
Y.~He \textit{et al.},
``TransMLA: Multi-head latent attention is all you need,''
\textit{arXiv:2502.07864}, 2025.

\bibitem{xkv}
C.-C.~Chang \textit{et al.},
``xKV: Cross-layer SVD for KV-cache compression,''
\textit{arXiv:2503.18893}, 2025.

\bibitem{longbench}
Y.~Bai \textit{et al.},
``LongBench: A bilingual, multitask benchmark for long context
understanding,''
\textit{arXiv:2308.14508}, 2023.

\end{thebibliography}

\end{document}